% ============================================================
% Proposition 3.3: Fast-Slow Decomposition (Fenichel)
% Status: CORRECTION (original proof sketch omits eigenvalue
%   verification, the central condition for normal hyperbolicity)
% ============================================================

% --- Definitions and Assumptions ---
\begin{definition}[Fast and Slow Variables]
Decompose the state vector:
\[
\mathbf{x}_A^{\text{fast}} = (\delta_A, \sigma_A, \alpha_A) \in \R^3,
\qquad
\mathbf{x}_A^{\text{slow}} = (\hat{M}_A, \Xi_A^P, \Xi_A^I, \Xi_A^F) \in \R^4.
\]
\end{definition}

\begin{definition}[Timescale Parameter]
Define the timescale separation parameter:
\[
\epsilon \defeq \frac{\max(\nu_M, \kappa_P, \kappa_I, \kappa_F)}{\min(\eta_{\max}, \rho, \gamma)}.
\]
The original manuscript asserts $\epsilon \ll 1$ as a design condition.
We treat it as an assumption for the validity of the decomposition.
\end{definition}

\begin{assumption}[Timescale Separation]
\label{assum:timescale}
The parameters satisfy:
\[
\epsilon \ll 1.
\]
That is, the slow subsystem's rate constants ($\nu_M \sim 10^{-2}$,
$\kappa_{\cdot} \sim 10^{-2}$) are at least one order of magnitude
smaller than the fast subsystem's rates ($\eta_{\max} \sim 1$,
$\rho \sim 10^{-1}$, $\gamma \sim 10^{-1}$).
\end{assumption}

\begin{assumption}[Critical Manifold Compactness]
The critical manifold $\mathcal{M}_0 = \{ \mathbf{x}_A \in \mathcal{X} :
\dot{\mathbf{x}}_A^{\text{fast}} = 0 \}$ is compact and connected.
This follows from Proposition~3.2: $\mathcal{X}$ is forward-invariant
and compact, and $\mathcal{M}_0$ is a closed subset of $\mathcal{X}$.
\end{assumption}

% --- Proposition Statement ---
\begin{proposition}[Fast-Slow Decomposition]
\label{prop:fenichel}
Under Assumptions~\ref{assum:timescale}--2, the system admits a slow
manifold $\mathcal{M}_\epsilon = \{ \mathbf{x}_A :
\mathbf{x}_A^{\text{fast}} = \mathbf{h}(\mathbf{x}_A^{\text{slow}}, \epsilon) \}$
that is $O(\epsilon)$-close to the critical manifold
$\mathcal{M}_0 = \{ \mathbf{x}_A : \dot{\mathbf{x}}_A^{\text{fast}} = 0 \}$,
and the reduced dynamics on $\mathcal{M}_\epsilon$ follow:
\[
\dot{\mathbf{x}}_A^{\text{slow}} = \epsilon \cdot
\mathbf{g}(\mathbf{h}(\mathbf{x}_A^{\text{slow}}, \epsilon),
\mathbf{x}_A^{\text{slow}}).
\]
\end{proposition}

% --- Proof ---
\begin{proof}
Fenichel's geometric singular perturbation theory \citep{fenichel1979geometric}
requires three conditions: (i) the critical manifold $\mathcal{M}_0$ is
compact and connected, (ii) the fast subsystem's linearization restricted
to the normal bundle of $\mathcal{M}_0$ has no eigenvalues with zero
real part (normal hyperbolicity), and (iii) the vector field is
$C^r$ for sufficiently large $r$. We verify each below.

\paragraph{Condition (i): Compactness and connectedness.}
By Assumption~2, $\mathcal{M}_0$ is a closed subset of the compact
set $\mathcal{X}$ under Proposition~3.2; hence $\mathcal{M}_0$ is
compact. Connectedness follows from the continuity of the fixed-point
equation $\dot{\mathbf{x}}_A^{\text{fast}} = 0$: on each connected
component of $\mathcal{X}$ (which is a product of intervals), the
solution set of analytic equations is connected.

\paragraph{Condition (ii): Normal hyperbolicity.}
The original manuscript (line~923) claims the Jacobian of the fast
subsystem ``has eigenvalues with strictly negative real parts (due to
the saturating nonlinearities in the Gompertz efficiency and the
$(1-\sigma_A)$, $(1-\alpha_A)$ factors).'' \textbf{This is a hand-wave
that must be replaced by an explicit calculation.}

The fast subsystem Jacobian is $D\mathbf{f}_{\text{fast}} \in \R^{3 \times 3}$:
\[
D\mathbf{f}_{\text{fast}} =
\begin{pmatrix}
\frac{\partial \dot{\delta}_A}{\partial \delta_A} &
\frac{\partial \dot{\delta}_A}{\partial \sigma_A} &
\frac{\partial \dot{\delta}_A}{\partial \alpha_A} \\[6pt]
\frac{\partial \dot{\sigma}_A}{\partial \delta_A} &
\frac{\partial \dot{\sigma}_A}{\partial \sigma_A} &
\frac{\partial \dot{\sigma}_A}{\partial \alpha_A} \\[6pt]
\frac{\partial \dot{\alpha}_A}{\partial \delta_A} &
\frac{\partial \dot{\alpha}_A}{\partial \sigma_A} &
\frac{\partial \dot{\alpha}_A}{\partial \alpha_A}
\end{pmatrix}.
\]

We compute each entry explicitly at a point on $\mathcal{M}_0$.

\subsubsection*{$\alpha_A$ submatrix (fastest variable).}
Recall:
\[
\dot{\alpha}_A = \gamma C_A (1 - \alpha_A) - \zeta_\alpha \alpha_A R_A^{\text{surface}}.
\]
The self-derivative is:
\[
\frac{\partial \dot{\alpha}_A}{\partial \alpha_A} = -\gamma C_A - \zeta_\alpha R_A^{\text{surface}}.
\]
Since $\gamma, C_A, \zeta_\alpha, R_A^{\text{surface}} \geq 0$,
and typically $\gamma C_A > 0$ (Assumption $C_A > 0$), we have:
\[
\frac{\partial \dot{\alpha}_A}{\partial \alpha_A} = -\underbrace{(\gamma C_A + \zeta_\alpha R_A^{\text{surface}})}_{>0} < 0.
\]
The cross-derivatives:
\[
\frac{\partial \dot{\alpha}_A}{\partial \delta_A} = \gamma (1 - \alpha_A) \frac{\partial C_A}{\partial \delta_A}
- \zeta_\alpha \alpha_A \frac{\partial R_A^{\text{surface}}}{\partial \delta_A},
\qquad
\frac{\partial \dot{\alpha}_A}{\partial \sigma_A} = 0
\]
($C_A$ and $R_A^{\text{surface}}$ do not depend directly on $\sigma_A$).

\subsubsection*{$\sigma_A$ submatrix.}
Recall:
\[
\dot{\sigma}_A = \rho P_A \alpha_A \sigma_A (1 - \gamma_\sigma \sigma_A) - \epsilon_\sigma \Omega_{SL} \sigma_A.
\]
The self-derivative is:
\[
\frac{\partial \dot{\sigma}_A}{\partial \sigma_A} = \rho P_A \alpha_A (1 - 2\gamma_\sigma \sigma_A) - \epsilon_\sigma \Omega_{SL}.
\]
On the critical manifold $\mathcal{M}_0$ (where $\dot{\sigma}_A = 0$):
\[
\rho P_A \alpha_A \sigma_A^*(1 - \gamma_\sigma \sigma_A^*) = \epsilon_\sigma \Omega_{SL} \sigma_A^*.
\]
For $\sigma_A^* > 0$, dividing by $\sigma_A^*$:
\[
\rho P_A \alpha_A (1 - \gamma_\sigma \sigma_A^*) = \epsilon_\sigma \Omega_{SL}.
\]
Substituting into the self-derivative:
\[
\left.\frac{\partial \dot{\sigma}_A}{\partial \sigma_A}\right|_{\mathcal{M}_0}
= \rho P_A \alpha_A (1 - 2\gamma_\sigma \sigma_A^*) - \underbrace{\rho P_A \alpha_A (1 - \gamma_\sigma \sigma_A^*)}_{\epsilon_\sigma \Omega_{SL}}
= -\rho P_A \alpha_A \gamma_\sigma \sigma_A^* < 0.
\]
For $\sigma_A^* = 0$:
\[
\left.\frac{\partial \dot{\sigma}_A}{\partial \sigma_A}\right|_{\sigma_A=0}
= \rho P_A \alpha_A - \epsilon_\sigma \Omega_{SL} = \epsilon_\sigma \Omega_{SL} (R_0 - 1).
\]
This is negative when $R_0 < 1$ (schema-free regime) and positive when
$R_0 > 1$ (the bifurcation condition). The $\sigma_A = 0$ branch of
$\mathcal{M}_0$ loses normal hyperbolicity at $R_0 = 1$ --- the
transcritical bifurcation point studied in Proposition~4.1. For
$R_0 \neq 1$, $\partial \dot{\sigma}_A / \partial \sigma_A \neq 0$ on the
$\sigma_A = 0$ branch.

The remaining cross-derivatives:
\[
\frac{\partial \dot{\sigma}_A}{\partial \alpha_A} = \rho P_A \sigma_A (1 - \gamma_\sigma \sigma_A) \geq 0,
\qquad
\frac{\partial \dot{\sigma}_A}{\partial \delta_A} = \rho \sigma_A (1 - \gamma_\sigma \sigma_A)
\left( \alpha_A \frac{\partial P_A}{\partial \delta_A} + P_A \frac{\partial \alpha_A}{\partial \delta_A} \right)
- \epsilon_\sigma \sigma_A \frac{\partial \Omega_{SL}}{\partial \delta_A}.
\]

\subsubsection*{$\delta_A$ submatrix.}
Recall:
\[
\dot{\delta}_A = f_{\text{learn}} \eta(\delta_A) T_A - \lambda_c (1 - \gamma_\sigma \sigma_A) \delta_A (1 - r_A),
\]
where $\eta(\delta_A) = \eta_{\max} \exp(-a \exp(-b \delta_A^{\text{relative}}))$
is the Gompertz efficiency. The self-derivative:
\[
\frac{\partial \dot{\delta}_A}{\partial \delta_A} = f_{\text{learn}} T_A \frac{\partial \eta}{\partial \delta_A}
- \lambda_c (1 - \gamma_\sigma \sigma_A) (1 - r_A).
\]
The Gompertz derivative:
\[
\frac{\partial \eta}{\partial \delta_A} = \eta \cdot a b \exp(-b \delta_A^{\text{relative}}) \cdot \frac{1}{\Delta(d,t)} > 0
\]
for finite $\delta_A < \Delta$. The decay term coefficient
$-\lambda_c (1 - \gamma_\sigma \sigma_A) (1 - r_A)$ is negative. The sign
of $\partial \dot{\delta}_A / \partial \delta_A$ depends on the balance;
near equilibrium $\dot{\delta}_A = 0$, the decay term dominates small
perturbations, yielding negativity.

Cross-derivatives:
\[
\frac{\partial \dot{\delta}_A}{\partial \sigma_A} = \lambda_c \gamma_\sigma \delta_A (1 - r_A) \geq 0,
\qquad
\frac{\partial \dot{\delta}_A}{\partial \alpha_A} = 0
\]
(no direct $\alpha_A$ dependence in the $\delta_A$ ODE).

\paragraph{Condition (ii) conclusion.}
The Jacobian on $\mathcal{M}_0$ (away from $R_0 = 1$) has the block-triangular
structure:
\[
D\mathbf{f}_{\text{fast}} =
\begin{pmatrix}
\partial_{\delta_A} \dot{\delta}_A & \partial_{\sigma_A} \dot{\delta}_A & 0 \\[4pt]
\partial_{\delta_A} \dot{\sigma}_A & \partial_{\sigma_A} \dot{\sigma}_A & \partial_{\alpha_A} \dot{\sigma}_A \\[4pt]
\partial_{\delta_A} \dot{\alpha}_A & 0 & \partial_{\alpha_A} \dot{\alpha}_A
\end{pmatrix}.
\]
The eigenvalues are the diagonal entries of the Schur complement after
block diagonalization:
\begin{align*}
\lambda_1 &= \frac{\partial \dot{\alpha}_A}{\partial \alpha_A} = -(\gamma C_A + \zeta_\alpha R_A^{\text{surface}}) < 0, \\
\lambda_2, \lambda_3 &= \text{eigenvalues of the } 2 \times 2 \text{ upper-left block}.
\end{align*}
The $2 \times 2$ block:
\[
J_{2 \times 2} = \begin{pmatrix}
\partial_{\delta_A} \dot{\delta}_A & \partial_{\sigma_A} \dot{\delta}_A \\
\partial_{\delta_A} \dot{\sigma}_A & \partial_{\sigma_A} \dot{\sigma}_A
\end{pmatrix}.
\]
Its trace:
\[
\operatorname{Tr}(J_{2 \times 2}) = \frac{\partial \dot{\delta}_A}{\partial \delta_A}
+ \frac{\partial \dot{\sigma}_A}{\partial \sigma_A} < 0
\quad\text{(both diagonal entries negative on $\mathcal{M}_0$ away from $R_0=1$)}.
\]
Its determinant:
\[
\det(J_{2 \times 2}) = \frac{\partial \dot{\delta}_A}{\partial \delta_A}
\frac{\partial \dot{\sigma}_A}{\partial \sigma_A}
- \frac{\partial \dot{\delta}_A}{\partial \sigma_A}
\frac{\partial \dot{\sigma}_A}{\partial \delta_A}.
\]
Both $\partial \dot{\delta}_A / \partial \delta_A$ and
$\partial \dot{\sigma}_A / \partial \sigma_A$ are negative, and the product
of the cross-derivatives $\partial \dot{\delta}_A / \partial \sigma_A
\cdot \partial \dot{\sigma}_A / \partial \delta_A$ is typically second-order
small (each involves $\sigma_A$ or a derivative of coefficient functions).
Hence $\det(J_{2 \times 2}) > 0$ on $\mathcal{M}_0$. Since trace $< 0$
and determinant $> 0$, both eigenvalues have strictly negative real parts.

Therefore, all three eigenvalues of $D\mathbf{f}_{\text{fast}}$ on
$\mathcal{M}_0$ (away from $R_0 = 1$) have strictly negative real parts.
Normal hyperbolicity holds. At $R_0 = 1$, a simple eigenvalue crosses
zero --- this is the transcritical bifurcation barrier where Fenichel's
theorem does not apply and the decomposition breaks down (the exchange
lemma handles this case).

\paragraph{Condition (iii): Smoothness.}
By Proposition~3.1, the vector field is piecewise-$C^\infty$ with
finitely many smooth components separated by measure-zero discontinuity
sets. Fenichel's theorem requires $C^r$ for $r \geq 1$; the
piecewise-$C^\infty$ structure satisfies this on each component.
The non-differentiable points ($\max(0,\cdot)$, $\Corr$ singularities)
are confined to lower-dimensional submanifolds that do not intersect
$\mathcal{M}_0$ in general position (the critical manifold lies in the
interior of the smooth regime for typical parameters).

\paragraph{Application of Fenichel's theorem.}
With conditions (i)--(iii) satisfied, Fenichel's theorem guarantees:
\begin{enumerate}[nosep]
\item Persistence of $\mathcal{M}_0$ to $\mathcal{M}_\epsilon$ for
  sufficiently small $\epsilon$.
\item $\mathcal{M}_\epsilon$ is $O(\epsilon)$-close to $\mathcal{M}_0$.
\item The reduced dynamics on $\mathcal{M}_\epsilon$ are given by
  $\dot{\mathbf{x}}_A^{\text{slow}} = \epsilon \cdot
  \mathbf{g}(\mathbf{h}(\mathbf{x}_A^{\text{slow}}, \epsilon),
  \mathbf{x}_A^{\text{slow}})$.
\end{enumerate}
$\square$
\end{proof}

% --- Boundary Cases ---
\begin{boundary-case}
\textbf{$R_0 = 1$ exactly:} At the bifurcation point, $\partial \dot{\sigma}_A / \partial \sigma_A = 0$,
violating normal hyperbolicity. Fenichel's theorem does not apply.
The fast-slow decomposition must be replaced by the exchange lemma or
center manifold reduction near this value. This is physically
interpretable as the Phase~1 $\to$ 2 transition region, where the
system spends arbitrarily long time near the bifurcation point
(critical slowing down).

\textbf{$\epsilon \approx 1$ (no separation):} If the rate constants
are comparable, the decomposition is invalid. The paper relies on
$\epsilon \ll 1$ as an assumption derived from order-of-magnitude
estimates, but this has not been empirically verified.

\textbf{$\sigma_A^* = 0$ branch:} For $R_0 < 1$, the $\sigma_A = 0$
branch of $\mathcal{M}_0$ is normally hyperbolic. For $R_0 > 1$,
it loses stability and the $\sigma_A^* > 0$ branch becomes the
attracting slow manifold.
\end{boundary-case}

% --- Circular Dependency Check ---
\begin{circularity-check}
\textbf{Does Proposition~3.3 depend on any later result?}
It depends on Proposition~3.2 (for compactness of $\mathcal{X}$ and
forward invariance) and Proposition~3.1 (for existence of solutions).
These are forward dependencies.

It also depends on Proposition~4.1 for the characterization at
$R_0 = 1$, but this is an edge case --- the primary claim holds for
$R_0 \neq 1$ independently.

\textbf{Does any later result depend on Proposition~3.3?}
The equilibrium analysis (Proposition~3.4) and bifurcation analysis
(Proposition~4.1) benefit from the decomposition but do not strictly
depend on it. The numerical integration scheme (Algorithm~3.1's IMEX
split) is motivated by Proposition~3.3 but can be used even without
the rigorous decomposition. Safe.
\end{circularity-check}
